\scanpage{15}\mspage{13}
Ceci dit, on a ceci (revenant à la notation $C_{\infty}$ ambivalente) : \inserted{a)} les inclusions \unsure{ple|pl.} fid.
\diagram{\begin{tikzcd}[ampersand replacement=\&]
\unsure{\text{a)}}\quad C_{\infty}^{\deleted{\unsure{-|+}}}(a,b) \arrow[r,hook] \& C_{\infty}^{\deleted{\unsure{-|+}}}(a,b) \& C_{\infty}^{+}(a,b) \arrow[l,hook']
\end{tikzcd}}
sont des iso dans $\mathrm{Ind}\,\mathrm{Cat}\,\mathrm{hot}$ resp. dans $\mathrm{Hot}$, et

b) On a bel et bien \unsure{une} composition des chemins
\[
C_{\infty}^{-}(a,b)\times C_{\infty}^{+}(b,c)\longrightarrow C_{\infty}(a,c).
\]
L'\uline{associativité} dans le cas \unsure{du|des} $\mathrm{CH}$ n'offre d'ailleurs pas de difficulté, à cause de la possibilité de se \deleted{\unsure{laisser}} \unsure{ramener|revenir} aux
\[
C\deleted{\unsure{H}}_{\unsure{m|m_1}}(a,b)\times C\deleted{\unsure{H}}_n(b,c)\times C_p(c,d)\rightrightarrows C_N(a,d)
\]
pour $N$ grand, et prouver l'égalité \deleted{\unsure{dans C}} comme flèches dans $\mathrm{Cat}\ \mathrm{hot}$ -- ce qui résultera des lemmes d'homotopie sur les intervalles dans $\mathrm{Cat}$. Dans le cas des $\mathrm{Ch}_{\infty}$ par contre je n'ai pas trouvé de démonstration. (Faute seulement d'avoir essayé sérieusement de m'y mettre ? Je n'ai pas trop envie, vu qu'il se pourrait que $\mathrm{Hot}\unsure{\text{---}|\to}\mathrm{Hot}_0^{\wedge}$ soit pleinement fidèle, auquel cas le travail devient \unsure{inutile} \ldots)
% archive sheet number: 14

